% =========================================================================
% CONECTIVIDADE, ROTAS DISJUNTAS E DECOMPOSIÇÃO
% =========================================================================
\chapter{Conectividade, Rotas Disjuntas e Decomposição}\label{sec:conectividade_roteamento}

A determinação do número de rotas independentes entre dois vértices $s$ e $t$ em um digrafo reduz-se diretamente ao cálculo de fluxo máximo sob duas formulações clássicas: caminhos disjuntos por arcos e por vértices.


\section{Caminhos Disjuntos por Arcos e por Vértices}\label{sec:caminhos_disjuntos}

No Problema dos Caminhos Disjuntos por Arcos (\textit{Edge-Disjoint Paths}), busca-se o número máximo de caminhos de $s$ a $t$ que não compartilham arcos.

A redução a fluxo máximo consiste em atribuir capacidade unitária a todos os arcos ($c(u,v) = 1, \forall (u,v) \in A$). Pelo Teorema da Integralidade (Teorema~\ref{teo:integralidade}), o fluxo máximo satisfaz $f(u,v) \in \{0, 1\}$. O valor ótimo $|f^*|$ é igual ao número máximo de caminhos disjuntos por arcos entre $s$ e $t$~\cite{cook1998combinatorial}.

No Problema dos Caminhos Disjuntos por Vértices (\textit{Vertex-Disjoint Paths}), os caminhos não podem compartilhar vértices intermediários. Como as restrições de capacidade em redes de fluxo aplicam-se a arcos, utiliza-se a técnica de \textbf{Divisão de Vértices} (\textit{Vertex Splitting})~\cite{ahuja1993network}.

A técnica transforma o digrafo original $D = (V, A)$ em uma rede estendida $D' = (V', A')$:
\begin{enumerate}
    \item Cada vértice $v \in V \setminus \{s, t\}$ é dividido em dois vértices: $v_{in}$ (entrada) e $v_{out}$ (saída).
    \item Insere-se um arco interno $(v_{in}, v_{out})$ com capacidade unitária $c(v_{in}, v_{out}) = 1$.
    \item Arcos incidentes a $v$ em $D$ são direcionados para $v_{in}$.
    \item Arcos emergentes de $v$ em $D$ passam a sair de $v_{out}$ com capacidade 1 (ou $\infty$).
\end{enumerate}

O arco interno $(v_{in}, v_{out})$ limita o trânsito por $v$ a no máximo 1 unidade de fluxo, garantindo que nenhum vértice intermediário participe de mais de um caminho. O valor do fluxo máximo em $D'$ equivale ao número máximo de caminhos disjuntos por vértices entre $s$ e $t$ em $D$~\cite{korte2018combinatorial}.

\begin{figure}[!ht]
\centering
\begin{tikzpicture}[thick]

\begin{scope}[xshift=0cm]
  \tikzsubcaption{(1.5,2.8)}{(a) Original $D$}
  \node[source node] (s) at (0,1.5) {$s$};
  \node[flow node] (v) at (1.5,1.5) {$v$};
  \node[sink node] (t) at (3,1.5) {$t$};
  \draw[flow edge] (s) -- (v); \draw[flow edge] (v) -- (t);
\end{scope}

\node at (4.5,1.5) {\Huge $\Rightarrow$};

\begin{scope}[xshift=6cm]
  \tikzsubcaption{(3.0,2.8)}{(b) Vertex Splitting: $D'$}
  \node[source node] (s2) at (0,1.5) {$s$};
  \node[flow node,fill=blue!10] (vin) at (2.2,1.5) {$v_{in}$};
  \node[flow node,fill=blue!20] (vout) at (4.2,1.5) {$v_{out}$};
  \node[sink node] (t2) at (6.4,1.5) {$t$};
  \draw[flow edge] (s2) -- node[above,font=\footnotesize]{$\infty$} (vin);
  \draw[sat edge] (vin) -- node[above,font=\footnotesize,red]{$c=1$} (vout);
  \draw[flow edge] (vout) -- node[above,font=\footnotesize]{$\infty$} (t2);
  \node[font=\footnotesize,below] at (3.2,0.7) {arco interno (gargalo)};
\end{scope}
\end{tikzpicture}
\caption{Vertex Splitting: cada vértice intermediário $v$ é dividido em $v_{in}$ e $v_{out}$ com arco interno de capacidade~1, limitando o fluxo pelo vértice.}
\label{fig:divisao_vertices}
\end{figure}


\section{O Teorema de Menger e as Dualidades da Conectividade}

O Teorema de Menger (1927) estabelece a equivalência min-max entre caminhos disjuntos e cortes de desconexão em digrafos, correspondendo à aplicação do Teorema Max-Flow Min-Cut sobre capacidades unitárias~\cite{ahuja1993network, cook1998combinatorial}.

\begin{teorema}[Menger --- Versão por Arcos, 1927][teo:menger_arcos]
Sejam $D = (V, A)$ um digrafo e $s, t \in V$ distintos e não adjacentes. O número máximo de caminhos direcionados de $s$ a $t$ disjuntos por arcos é igual ao número mínimo de arcos cuja remoção desconecta todo caminho de $s$ a $t$:
\begin{equation}
    \max\{\text{caminhos $st$-disjuntos por arcos}\} = \min\{\text{arcos cuja remoção desconecta $s$ de $t$}\}
\end{equation}
\end{teorema}

\begin{demonstracao}
Defina $c(u,v) = 1$ para todo $(u,v) \in A$. Pelo Teorema da Integralidade (Teorema~\ref{teo:integralidade}), o fluxo máximo é inteiro ($f(u,v) \in \{0, 1\}$). O valor $|f^*|$ é igual ao número de caminhos disjuntos por arcos. Pelo Teorema~\ref{teo:mfmc}, $|f^*| = \ccap(S,T)$ para algum corte mínimo $(S,T)$. Como as capacidades são unitárias, $\ccap(S,T)$ coincide com o número mínimo de arcos que desconectam $s$ de $t$.
\end{demonstracao}

\begin{teorema}[Menger --- Versão por Vértices, 1927][teo:menger_vertices]
Sejam $D = (V, A)$ um digrafo e $s, t \in V$ distintos e não adjacentes. O número máximo de caminhos direcionados de $s$ a $t$ internamente disjuntos por vértices é igual ao número mínimo de vértices cuja remoção desconecta todo caminho de $s$ a $t$:
\begin{equation}
    \max\{\text{caminhos $st$-disjuntos por vértices}\} = \min\{\text{vértices cuja remoção desconecta $s$ de $t$}\}
\end{equation}
\end{teorema}

\begin{demonstracao}
Aplica-se \textit{Vertex Splitting} sobre $D$ para obter $D'$ com capacidades unitárias nos arcos internos $(v_{in}, v_{out})$ e infinitas nos demais arcos. Pela versão por arcos aplicada a $D'$, o corte mínimo é composto exclusivamente por arcos internos, cuja quantidade equivale ao número mínimo de vértices de $D$ que desconectam $s$ de $t$.
\end{demonstracao}


\section{Transformações Estruturais de Redes de Fluxo}

\subsection{Múltiplas Fontes e Múltiplos Sumidouros}

Redes com múltiplas fontes $s_1, \ldots, s_k$ e sumidouros $t_1, \ldots, t_m$ reduzem-se à formulação padrão $s$-$t$ introduzindo uma \textbf{superfonte} $s^*$ e um \textbf{supersumidouro} $t^*$~\cite{ahuja1993network}:
\begin{enumerate}
    \item Para cada fonte $s_i$, adiciona-se o arco $(s^*, s_i)$ com capacidade igual à oferta líquida de $s_i$ (ou $\infty$).
    \item Para cada sumidouro $t_j$, adiciona-se o arco $(t_j, t^*)$ com capacidade igual à demanda líquida de $t_j$ (ou $\infty$).
\end{enumerate}

\begin{figure}[!ht]
\centering
\begin{tikzpicture}[thick]
    \node[source node] (ss) at (0, 0) {$s^*$};
    \node[flow node, fill=green!10] (s1) at (2, 1.5) {$s_1$};
    \node[flow node, fill=green!10] (s2) at (2, -1.5) {$s_2$};

    \node[flow node] (v1) at (4, 1.5) {$v_1$};
    \node[flow node] (v2) at (4, -1.5) {$v_2$};
    \node[flow node] (v3) at (5.5, 0) {$v_3$};

    \node[flow node, fill=red!10] (t1) at (7, 1.5) {$t_1$};
    \node[flow node, fill=red!10] (t2) at (7, -1.5) {$t_2$};

    \node[sink node] (ts) at (9, 0) {$t^*$};

    \draw[->, dashed, blue, line width=1.2pt] (ss) -- node[above left, font=\footnotesize] {$\infty$} (s1);
    \draw[->, dashed, blue, line width=1.2pt] (ss) -- node[below left, font=\footnotesize] {$\infty$} (s2);

    \draw[flow edge] (s1) -- (v1);
    \draw[flow edge] (s1) -- (v3);
    \draw[flow edge] (s2) -- (v2);
    \draw[flow edge] (s2) -- (v3);

    \draw[flow edge] (v1) -- (t1);
    \draw[flow edge] (v1) -- (v3);
    \draw[flow edge] (v3) -- (t1);
    \draw[flow edge] (v3) -- (t2);
    \draw[flow edge] (v2) -- (t2);

    \draw[->, dashed, red, line width=1.2pt] (t1) -- node[above right, font=\footnotesize] {$\infty$} (ts);
    \draw[->, dashed, red, line width=1.2pt] (t2) -- node[below right, font=\footnotesize] {$\infty$} (ts);
\end{tikzpicture}
\caption{Redução de múltiplas fontes e sumidouros via superfonte $s^*$ e supersumidouro $t^*$.}
\label{fig:multiplas_fontes}
\end{figure}

O fluxo máximo de $s^*$ a $t^*$ em $D'$ equivale ao fluxo máximo viável entre o conjunto de fontes e o conjunto de sumidouros em $D$.


\subsection{Arestas Bidirecionais e Grafos Não Direcionados}

Para aplicar algoritmos de fluxo em grafos não direcionados $G = (V, E)$, cada aresta $\{u,v\} \in E$ com capacidade $c$ é substituída por dois arcos antiparalelos direcionados $(u,v)$ e $(v,u)$, ambos com capacidade $c$~\cite{ahuja1993network}.

Na implementação, a chamada \lstinline{add_edge(u, v, c, c)} inicializa ambos os arcos com capacidade $c$, mantendo o emparelhamento por XOR no vetor de arcos.

\begin{figure}[!ht]
\centering
\begin{tikzpicture}[thick]
  \begin{scope}[shift={(0,0)}]
    \tikzsubcaption{(1.5,1.5)}{(a) Aresta Não-direcionada}
    \node[trans node] (u) at (0, 0) {$u$};
    \node[trans node] (v) at (3, 0) {$v$};
    \draw[-] (u) -- node[above, font=\footnotesize] {$c$} (v);
  \end{scope}

  \begin{scope}[shift={(6,0)}]
    \tikzsubcaption{(1.5,1.5)}{(b) Arcos Direcionados}
    \node[trans node] (u2) at (0, 0) {$u$};
    \node[trans node] (v2) at (3, 0) {$v$};
    \draw[->, transform canvas={yshift=1.5mm}] (u2) -- node[above, font=\footnotesize] {$c$} (v2);
    \draw[<-, transform canvas={yshift=-1.5mm}] (u2) -- node[below, font=\footnotesize] {$c$} (v2);
  \end{scope}
\end{tikzpicture}
\caption{Modelagem de aresta não direcionada por meio de dois arcos direcionados opostos com capacidade $c$.}
\label{fig:arcos_bidirecionais}
\end{figure}


\section{Decomposição de Fluxo}\label{sec:decomposicao_cap}

\subsection{Decomposição em Caminhos e Ciclos}

\begin{teorema}[Decomposição de Fluxo][teo:decomposicao]
Todo fluxo viável $f$ em $D = (V, A)$ pode ser decomposto em uma família de no máximo $|A|$ trajetórias simples (caminhos de $s$ a $t$ e ciclos direcionados), cada uma com fluxo estritamente positivo. A soma dos fluxos dos caminhos iguala $|f|$, e o fluxo em cada arco é a soma dos fluxos das trajetórias que o contêm.
\end{teorema}

\intuicao{Decomposição Linear} O fluxo em cada arco resulta da sobreposição de caminhos de $s$ a $t$ e de ciclos de circulação. A cada passo indutivo, extrai-se uma trajetória com gargalo $\delta > 0$, zerando o fluxo em pelo menos um arco e garantindo término em no máximo $|A|$ etapas.

\begin{demonstracao}
Por indução no número de arcos com fluxo positivo $|A_f^+| = |\{(u,v) \in A \mid f(u,v) > 0\}|$:

\medskip
\noindent \textit{Base ($|A_f^+| = 0$):} Fluxo nulo, decomposição vazia.

\medskip
\noindent \textit{Passo indutivo:} Suponha válido para menos de $|A_f^+|$ arcos positivos. Se $|f| > 0$, parte-se de $s$ seguindo arcos com $f > 0$. Pela conservação, a trajetória alcança $t$ (caminho simples $P$) ou fecha um ciclo simples $C$. Se $|f| = 0$ com $|A_f^+| > 0$, a conservação em todos os vértices garante a existência de um ciclo simples $C$.

Seja $\mathcal{T}$ a trajetória obtida (caminho ou ciclo) e $\delta = \min_{a \in \mathcal{T}} f(a) > 0$. Subtraindo $\delta$ ao longo de $\mathcal{T}$, obtém-se fluxo viável $f'$ com $|A_{f'}^+| \le |A_f^+| - 1$. Por indução, $f'$ admite decomposição em no máximo $|A| - 1$ elementos; somando $\mathcal{T}$ com valor $\delta$, obtém-se a decomposição de $f$ com no máximo $|A|$ trajetórias.
\end{demonstracao}

A decomposição de fluxo é utilizada na caracterização da otimalidade por ciclos negativos em Fluxo de Custo Mínimo (Seção~\ref{sec:condicoes_otimalidade}), onde a diferença entre dois fluxos de mesmo valor constitui uma circulação pura.

\begin{figure}[!ht]
\centering
\begin{tikzpicture}[thick]
  \begin{scope}[shift={(0,0)}]
    \node[source node, font=\footnotesize] (s) at (0, 1) {$s$};
    \node[flow node, font=\footnotesize] (a) at (2, 2) {$a$};
    \node[flow node, font=\footnotesize] (b) at (2, 0) {$b$};
    \node[sink node, font=\footnotesize] (t) at (4, 1) {$t$};

    \draw[flow edge] (s) -- node[edge lbl above] {$3$} (a);
    \draw[flow edge] (s) -- node[edge lbl below] {$2$} (b);
    \draw[flow edge] (a) -- node[edge lbl right] {$1$} (b);
    \draw[flow edge] (a) -- node[edge lbl above] {$2$} (t);
    \draw[flow edge] (b) -- node[edge lbl below] {$3$} (t);
  \end{scope}

  \begin{scope}[shift={(6,0.5)}]
    \node[right, font=\small, align=left] at (0, 1.5) {\textbf{Caminhos Decompostos}:};
    \node[right, font=\footnotesize, align=left] at (0, 0.8) {1) $s \rightarrow a \rightarrow t$ (fluxo 2)};
    \node[right, font=\footnotesize, align=left] at (0, 0.2) {2) $s \rightarrow b \rightarrow t$ (fluxo 2)};
    \node[right, font=\footnotesize, align=left] at (0, -0.4) {3) $s \rightarrow a \rightarrow b \rightarrow t$ (fluxo 1)};
  \end{scope}
\end{tikzpicture}
\caption{Exemplo de decomposição de fluxo: o fluxo total de 5 unidades é expresso como soma de três caminhos direcionados de $s$ a $t$.}
\label{fig:decomposicao_fluxo}
\end{figure}


\subsection{Reconstrução de Caminhos}

Em problemas com capacidades unitárias, a extração dos caminhos após a terminação do cálculo de fluxo máximo opera percorrendo os arcos com $f(u,v) = 1$:
\begin{enumerate}
    \item A partir de $s$, segue-se um arco com $f(u,v) = 1$ ainda não marcado.
    \item Marca-se o arco e avança-se até $t$, registrando o caminho.
    \item O processo repete-se $|f^*|$ vezes.
\end{enumerate}

Essa rotina é utilizada na resolução do problema \textit{Distinct Routes} para emitir as rotas disjuntas encontradas.


